% !TEX recipe = LuaLaTeX+se
\documentclass[10pt,parskip,paper=216mm:162mm]{graduale} % set document class: manual at https://ctan.org/pkg/koma-script

\begin{document}

\subsection{Copyright notices}

For example, the commentary for the following Introit:
\begin{center}
\emph{In excelso throno, GR 257, Cf. Dan. 7: 9-10, 13-14, Is. 6: 1-3 V. Ps. 99, tr. RM V. APC}
\end{center}

\begin{itemize}
    \item\emph{In excelso throno} --- the name of the chant in the Graduale Romanum
    \item\emph{GR 257} found on page 257 of the 1974 Graduale Romanum
    \item\emph{Cf. Dan. 7: 9-10, 13-14, Is. 6: 1-3 V. Ps. 99} Antiphon from Dan. 7: 9-10, 13-14, Is. 6: 1-3, verse from Psalm 100[99] (See style guideline no. 2)
    \item\emph{tr. RM V. APC} Translation of the antiphon taken from The Roman Missal, translation of the verse taken from The Abbey Psalms and Canticles. % Since APC uses the Hebrew numbering for the Psalms, the translation is taken from \emph{Psalm 100} in APC even though the commentary says Psalm 99.
\end{itemize}

\subsubsection{Citations for chant sources}
\begin{itemize}
    \item[GR] - Graduale Romanum 1974, Solesmes.
    \item[GR61] - Graduale Romanum 1961, Solesmes.
    \item[GR09] - Graduale Romanum 1909, Vatican edition.
    \item[OR] - Offertoriale 1935, Karl Ott.
    \item[AM] - Antiphonale Monasticum 1935, Solesmes.
\end{itemize}

\subsubsection{Citations for text sources}
\begin{itemize}
    \item[LFM] - Lectionary for Mass © 1969, 1981, 1997, International Commission on English in the Liturgy Corporation (ICEL);
    \item[RM] - The Roman Missal © 2010, ICEL. 
    \item[APC] - The Abbey Psalms and Canticles, 
    \item[GM] - Gregorian Missal, 1990, Solesmes.
    \item[NAB] - New American Bible,
    \item[LH] - The Liturgy of the Hours © 1973, 1974, 1975, ICEL.
\end{itemize}

\subsection{Foreword}

\subsubsection{What makes this book special}
I want to start off by giving a massive acknowledgement to Fr. Weber for
his book, \emph{The Proper of the Mass for Sundays and Solemnities}.

The \emph{Graduale Novi Mundi} was written to provide another option for
those who want to provide beautiful English chants that are more faithful
to those in the \emph{Graduale Romanum}. This book will build on what
Fr. Weber started by, on the one hand, providing translations of those
Introits and Communion chants whose texts are not present in the Missal,
and on the other, 


\subsubsection{When the Graduale and Missal have different texts}
If you've been to a daily Mass without singing, you may have noticed the
extra texts recited by the priest or the faithful at the
procession and as the priest receives Communion

From the very beginning of the Church, Mass has largely been chanted in
full. This included chants that would be sung to complement other liturgical
acts, such as

In medieval times, as the faithful desired to go to Mass more often, 

In \emph{The Proper of the Mass}, Fr. Weber has chosen to retain only the texts
given in the altar Missal for the Introit and Communion antiphons, and has set these
to various chants.

From the GIRM (United States edition), paragraph 48 provides (emphasis added):
\begin{center}
    \itshape
    In the Dioceses of the United States of America,
    there are four options for the Entrance Chant: \lbrack the first being \rbrack
    \textbf{the antiphon from the Missal \underline{or} the antiphon with its Psalm from the Graduale Romanum},
    as set to music there or in another setting\lbrack.\rbrack
\end{center}
The instructions for the Offertory and Communion chants follow similarly.
Unlike the Responsorial Psalm and Alleluia, which ordinarily \emph{do}
replace the Gradual and Alleluia in the \emph{Graduale},\footnote{In the United States, it is 
permissible to sing them, but it is not clear if it is permissible to translate them directly to the vernacular.}
the Introit, Offertory and Communion simply have two official liturgical texts and either may be chosen
when the Mass is sung.

\subsubsection{Style guidelines}
\begin{enumerate}
    \item The `\emph{j}' is used in all cases for the consonantal form of the `\emph{i}'.
        As such, the Offertory chant for the first week in Ordinary Time is cited as ``Jublate Deo omis''
        even though the 1974 Graduale Romanum prints ``Iubilate''.
    \item When a Psalm cited in the commentary field has no brackets, it is presumed to be
        Suptuagint numbering. However, in certain cases, both are given, in the form
        Hebrew[Septuagint],
        as in the Communion chant Feci judicium:
        \begin{center}
        ``Feci judicium, GR 529, Ps. 119[118]:121,101, tr. APC''
        \end{center}
\end{enumerate}

\end{document}